\section{Evaluation}
\label{chap:evaluation}

Die Evaluation verwendet nachweisbare Eigenschaften: stabile Hashes
projektverwalteter Dateien, ein schreibfreier Check, ein leerer zweiter Plan,
Wiederherstellung bei Fehlversuchen und die Konsistenz eines exportierten Payloads. Das
Design trennt \Implemented{} von \Verified{} und \Measured{}; die letzte Kategorie
enthält erst nach einem reproduzierbaren Lauf Versions-, Commit-, Runner- und
Toolchainangaben. Der portable Export ist als \Verified{}
(\evidence{EV-EXPORT-001}) registriert.

Abbildung~\ref{fig:export-boundary} verdeutlicht die Grenze zwischen dem exportierten
Toolingbestand und nicht portablen Projekt- oder Builddateien.

\CaseDiagram{export-boundary}{Exportgrenze: nur Tools und portable Dokumentation werden inventarisiert.}{fig:export-boundary}
\CaseTable{
\begin{tabularx}{0.94\textwidth}{@{}lX@{}}
\toprule
Beobachtung & Interpretation \\
\midrule
Kontextpfade liegen im Zielprojekt & Keine Quellrepository-Wurzel erforderlich \\
Zweiter Full-Fix plant nichts & Gewünschter Zustand und State stimmen überein \\
Fehler löst Rückgabe aus & Teilveröffentlichung wird nicht als Erfolg gewertet \\
\bottomrule
\end{tabularx}
}{Auswertungslogik ohne erfundene Messwerte.}{tab:evaluation-design}
